\section{Conclusion}
\label{sec:conclusion}

The delineation of the lesion margin is the principal difficulty of dermoscopic segmentation, and the aspect least well served by the region-overlap criteria with which methods are conventionally optimised and assessed. This work has investigated a targeted dual-stream refinement of the MAFF-Net fusion stage, directed at that aspect. The pretrained encoder, the convolutional decoder, and the auxiliary reconstruction branch are retained from the Original MAFF-Net, and the redesign is confined to the fusion pathway. The multi-scale features of the frozen backbone are processed along two parallel streams --- a spatial stream preserving regional and contextual structure, and a parameter-free Frequency Branch providing an explicit spectral route for boundary-sensitive content --- which a Feature Alignment Module reconciles and a Gated Fusion Module combines by a learned, spatially varying weighting. Training uses a region objective supplemented by auxiliary boundary-focused supervision.

Evaluated on three public dermoscopic benchmarks against four segmentation architectures under an identical protocol, all members of the MAFF-Net family exceeded every baseline on both families of measures across all three datasets. A leave-one-out ablation established that disabling any one of the three proposed components degrades all three boundary measures while leaving the region measures substantially unaffected --- the complete model recording the highest value on no region measure --- so their effect is confined to the fidelity of the predicted contour, and the results do not support, and this work does not advance, any claim that they improve region overlap. The proposed configuration was moreover not the strongest member of the MAFF-Net family, the restored variant recording the higher value on the majority of the measures.

Stated conservatively, the contribution of this work is threefold. First, evaluating the models with region-based and boundary-based measures in combination shows that the two families induce different rankings, so that region-overlap criteria alone are insufficient to characterise methods intended to act upon the lesion margin. Second, the proposed dual-stream model attains higher region-based and boundary-based scores than four standard convolutional and transformer baselines on all three benchmarks under an identical single-run protocol. Third, the dual-stream reformulation of the fusion pathway is investigated and the marginal role of its three components isolated through a controlled leave-one-out ablation, whose results confine the effect of these components to boundary fidelity rather than the extent of the predicted region. No comparison was drawn against results in the literature, and no claim of state-of-the-art performance is made.

The principal limitations indicate the directions in which the study should be extended. Every configuration was trained once, with a single experimental seed, and no statistical analysis was undertaken; repetition across several seeds, with appropriate interval estimates, is required before the smaller differences reported here can be regarded as established, and is the most immediate priority. The evaluation is confined to three benchmarks from a single imaging modality and annotation protocol, so assessment on further collections, and on other modalities and acquisition conditions, is needed to establish generality. Finally, the boundary-focused supervision terms were active in every configuration and were not ablated, so their individual contribution remains unmeasured; ablating them, and examining alternative formulations and their interaction with the architectural components, is a natural continuation of this study.
